\section{Method}

\subsection{Model}
\label{sec:method_model}

\textbf{Encoder.} One patch becomes one token (Fig.~\ref{fig:arch}(a)) through a \emph{single} linear projection; there is no codebook and no quantization. Channels are mixed at this first layer by design, since the factor of interest is sometimes defined by the relation between channels --- HAPT posture is determined by how gravity distributes over the three axes. Tokens then pass through a stack of \textbf{causal transformer layers} with rotary position encoding, and a final linear layer produces the embedding, whose leading coordinates form $\zslow$ and the remainder $\zmix$. The last operation is a \textbf{per-block LayerNorm}, applied to the two blocks separately rather than to the embedding as a whole, for the reason given in Section~\ref{sec:xcov}. The \textbf{predictor} is a small MLP that sees only the embedding at one time point together with the horizon, which enters as a learned expansion of $\log_2\Delta$; one predictor is conditioned on the horizon rather than one head per horizon, consistent with treating the horizon as a continuous axis.

\textbf{No modification is made to the attention operation}: the gate sits between encoder output and predictor input, the penalty in the loss, the per-block LayerNorm in the encoder's last line. This is a change of training procedure, not of architecture, so it attaches to any causal sequence encoder.

\begin{figure}[!tb]
\centering
\includegraphics[width=0.78\textwidth]{figs/fig_main.png}
\caption{Horizon-gated latent prediction. \textbf{(a)} A causal encoder maps a window to $z_t$, split into $\zslow$ (16 of 64 dimensions) and $\zmix$; a predictor conditioned on $\Delta$ matches a target read at one position inside $(t,\,t+\Delta]$. \textbf{(b, c)} The same pipeline at a short and a long horizon. \textbf{(d)} The gate, with the axis in units of $\tau$.}
\label{fig:arch}
\end{figure}

\textbf{Target.}
\begin{equation}
z^{*} = \mathrm{sg}\big[\,f_{\bar\theta}\big(x_{[t+\Delta-w,\;t+\Delta]}\big)_{-1}\big], \qquad \bar\theta \leftarrow m\bar\theta + (1-m)\theta
\label{eq:target}
\end{equation}

The prediction target is not the raw signal but the representation an EMA copy of the encoder reads from a short window inside the future interval, with the gradient stopped. Bounding that window keeps the target's receptive field inside $(t,\,t+\Delta]$, so it never re-contains the anchor's own past and the shortcut of matching the target from what is already known at the anchor is closed by construction. \textbf{The target never passes the gate and is always the full embedding}: at far horizons the model is asked to match the whole future embedding from the persistent block alone, and this asymmetry is the core of the mechanism.

\textbf{Loss.}
\begin{equation}
\mathcal{L} = \mathcal{L}_{\text{pred}} + \lambda\,\mathcal{L}_{\text{xcov}}, \qquad \mathcal{L}_{\text{pred}} \in \{\mathcal{L}_{1},\ \mathcal{L}_{\text{NCE}}\}
\label{eq:loss}
\end{equation}

The two forms share backbone, gate and penalty and differ only in the loss layer. \textbf{HGLP-Reg} uses the $L_1$ distance; $L_2$'s squared-error gradient is dominated by large residuals and buries a weak persistent signal, which made results swing across seeds. \textbf{HGLP-NCE} uses InfoNCE over cosine similarity with temperature 0.1, with one departure from the standard form: \textbf{samples from the same window are excluded from the negative pairs}, since they share the persistent factor and pushing them apart would teach the model to discard it.


\subsection{The Horizon Gate}
\label{sec:gate}

\begin{equation}
g(\Delta) = \sigma\!\Big(\frac{\tau-\Delta}{\kappa}\Big), \qquad \tilde z_t(\Delta) = \big[\,z_t^{\mathrm{per}}\,;\ g(\Delta)\,z_t^{\mathrm{mix}}\,\big]
\label{eq:gate}
\end{equation}

Immediately before the predictor, the rear block is multiplied by $g(\Delta)\in(0,1)$ (Fig.~\ref{fig:arch}(b)--(d)): for $\Delta\ll\tau$ the embedding passes unchanged, for $\Delta\gg\tau$ the rear block is suppressed. \emph{Dimensions are multiplied, not truncated}, so the predictor input is always $D$-dimensional; \emph{$\zslow$ is never blocked}; and the gate is \emph{a smooth curve, not a step}. The width $\kappa$ is a buffer rather than a convenience, since $\tau$ cannot be chosen accurately without labels, and a step gate was worse on both stems (Appendix~\ref{app:design}).

The one-sidedness is likewise a choice. Withholding $\zmix$ beyond $\tau$ should cost nothing, because the transient factor is already forgotten there. Withholding $\zslow$ below $\tau$ would cost a great deal: near horizons need the current persistent state as much as the fast detail, and they are where the gradients are dense, so the block would lose most of its training signal rather than gain protection. Appendix~\ref{app:design} tests the symmetric variant. Note also that only the \emph{availability} of $\zmix$ is constrained, never its content --- across the six settings it carries the persistent factor at \numMixPerLo--\numMixPerHi{} --- which is why we call it $\zmix$ and not ``$z_{\mathrm{fast}}$''.

\textbf{Gradient split.}
\begin{equation}
\frac{\partial\mathcal{L}}{\partial z_t^{\mathrm{per}}} = \frac{\partial\mathcal{L}}{\partial \tilde z_t^{\mathrm{per}}}, \qquad \frac{\partial\mathcal{L}}{\partial z_t^{\mathrm{mix}}} = g(\Delta)\,\frac{\partial\mathcal{L}}{\partial \tilde z_t^{\mathrm{mix}}}
\label{eq:grad}
\end{equation}

Because the encoder is shared across horizons, this one line splits the training signal the two blocks receive: $\zslow$ is updated at every horizon, whereas $\zmix$ is updated in proportion to $g(\Delta)$ and so effectively only at near horizons. Two consequences follow, and only one of them is guaranteed. \textbf{Availability is enforced}: a predictor cannot foresee when the persistent factor will next switch, so at long horizons it can lower the loss only by recording the current state, and the gate leaves $\zslow$ as the sole route for doing so; so whatever still contributes to prediction that far ahead must be readable there, and no predictor however large can restore information the gate has suppressed to zero. \textbf{Exclusion is not enforced}: $\zslow$ is still used at near horizons, where transient information lowers the loss, and no term of the objective pushes that information \emph{out of} $\zslow$ --- exclusion has to come from the penalty of Section~\ref{sec:xcov}. We therefore measure exclusion rather than assume it, and Section~\ref{sec:rotation} follows that measurement through.

\subsection{Pulling the Two Blocks Apart}
\label{sec:xcov}

\textbf{Cross-covariance penalty.}
\begin{equation}
\begin{split}
C &= \tfrac{1}{N-1}\,\bar Z^{\top}_{(1:\dslow)}\,\bar Z_{(\dslow+1:D)},\\[2pt]
\mathcal{L}_{\text{xcov}} &= \tfrac{1}{\dslow (D-\dslow)}\textstyle\sum_{a,b} C_{ab}^{2}
\end{split}
\label{eq:xcov}
\end{equation}

The penalty pushes the two blocks to be linearly unpredictable from each other and \textbf{acts on the embedding before the gate}. It travels with the \textbf{per-block LayerNorm}: a single LayerNorm over the whole embedding gives the two blocks one shared mean and variance, recreating exactly the coupling the penalty is trying to remove. \textbf{``Mechanism-free'' therefore means all three devices absent} --- gate, penalty and per-block LayerNorm.

\subsection{What We Measure}

Reading the persistent factor from $\zslow$ is not by itself evidence of separation, since rotational symmetry means it is read about as well from any 16 dimensions. We therefore report a \textbf{block $\times$ factor matrix} scoring both blocks on both factors, plus a single-number summary \textbf{SEP}, which asks three questions.

\begin{enumerate}
\item \textbf{Inclusion --- is the persistent factor in $\zslow$?} We score how well the persistent factor is predicted from $\zslow$ alone.
\item \textbf{Allocation --- is the transient factor in $\zmix$?} We score how well the transient proxy is recovered from $\zmix$ alone.
\item \textbf{Exclusion --- is the transient factor absent from $\zslow$?} We measure how well the transient proxy can be recovered from $\zslow$, and subtract that value from 1.
\end{enumerate}

\begin{equation}
\mathrm{SEP} = \underbrace{\big[s_{\text{per}}(z^{\mathrm{per}})\big]_0^1}_{\text{inclusion}} \cdot \underbrace{\big[s_{\text{tra}}(z^{\mathrm{mix}})\big]_0^1}_{\text{allocation}} \cdot \underbrace{\big[1-s_{\text{tra}}(z^{\mathrm{per}})\big]_0^1}_{\text{exclusion}}
\label{eq:sep}
\end{equation}

$s_{\text{per}}(\cdot)$ scores reading the \emph{persistent factor} from a subspace and $s_{\text{tra}}(\cdot)$ the \emph{transient proxy}; $[\cdot]_0^1$ clips into $[0,1]$. Classification scores already lie in this range, but $R^2$ goes negative when a prediction is worse than the mean, and the clipping acts only there, preventing a negative value from flipping the sign of the product. There is no normalization by any denominator anywhere.


\textbf{Probe protocol.} With the encoder frozen we fit only a linear classifier or regressor. Splits are \textbf{per group} (subject for HAPT, patient for PTB-XL, a front/back time split for the synthetic data), evaluation windows are \textbf{non-overlapping} to prevent leakage (asserted in the code), and only positions after the \textbf{minimum context length (64 patches)} are scored, since with less context the encoder lacks the evidence to read the persistent state. A \textbf{random subspace} of the same width is scored as a control, since rotational symmetry means the persistent factor is read about as well from any 16 dimensions.
